\section{Кратный интеграл}
\subsection{Определение кратного интеграла}
\begin{reminder}
	Интеграл Римана. \\
	Пусть $f$ ограничена на $[a,b]$, для неё есть разбиение $P:a=x_0<\dots<x_n=b$\\
	Вспомним диаметр разбиение, верхнюю и нижнюю суммы Дарбу и верхний и нижний интегралы:
	\[\tr(P)=\max_{1\le i\le n}\tr x_i,\ M_i=\sup_{[x_{i-1},x_i]}f(x),\ m_i=\inf_{[x_{i-1},x_i]}f(x)\]
	\[U(P,f)=\ns M_i\cdot\tr x_i,\ L(P,f)=\ns m_i\cdot\tr x_i\]
	\[\v I=\inf U(P,f),\ \uv I=\sup L(P,f)\]
	Если они равны, то $f$ называется интегрируемой по Риману на отрезке $[a,b]$, что также равно (по доказанной ранее теореме) пределу интегральных сумм:
	\[I=\v I=\uv I=\int_a^bf(x)dx=\lim_{\tr(P)\ra0}\ns f(t_i)\tr x_i\]
\end{reminder}
\begin{reminder}
	Вспомним определение измеримой функции.\\
	\textit{Выдержка из конспекта прошлого года:}
	\begin{definition}
		\textit{Лебеговым множеством} функции $f$, заданной на $E\subset\R^n,\ f:E\ra \R$, называется:
		\[E_f(a)=\{x\in E:f(x)<a\},\ a\in\R\]
		Числовая функция, заданная на измеримом по Лебегу множестве $E\subset\R^n$, называется \textit{измеримой}, если $(\fa a\in\R)\ E_f(a)$ измеримо.
	\end{definition}
\end{reminder}
\begin{definition}
	Пусть $E\subset\rn$ - измеримое по Лебегу (Жордану) множество конечной меры, $f:E\ra\R$ ограничена на $E$. Назовём $P$ \textit{разбиением} $E$, если: \[E=\bsc_{i=1}E_i\]
	где $E_i$ - измеримые по Лебегу (Жордану) множества.\\
	Положим:
	\[\tr(P)=\max_{1\le i\le N}\jm(E_i),\ M_i=\sup_{x\in E_i}f(x),\ m_i=\inf_{x\in E_i}f(x)\]
	\[U(P,f)=\ns M_i\cdot\jm(E_i),\ L(P,f)=\ns m_i\cdot\jm(E_i)\]
	Назовём верхним и нижним интегралом Лебега (Римана) следующие выражения соответственно:
	\[\v I=\inf U(P,f),\ \uv I=\sup L(P,f)\]
	Если верхний и нижний интегралы Лебега (Римана) равны, то $f$ называется \textit{интегрируемой по Лебегу на $E$}, обычно пишут:
	\[\v I=\uv I=\int_Ef(x)d\jm(x)\]
	В случае интегрируемости по Риману:
	\[\v I=\uv I=\int_Ef(x)dxdy\]
\end{definition} 
\begin{note}
	Далее будут использоваться сокращения:\vspace{-0.5cm}
	\[f\text{ интегрируема по Риману на множестве }E\lra f\in\mc R(E)\]
	\vspace{-1.23cm}
	\[f\text{ интегрируема по Лебегу на множестве }E\lra f\in\mc L(E)\]
\end{note}
\begin{theorem}[Критерий интегрируемости]
	\[f\in\mc R(E)\lra(\fa\eps>0)(\ex P)\ U(P,f)-L(P,f)<\eps\]
\end{theorem}
\begin{proof}
	Доказывается аналогично интегрируемости на отрезке. (а если конкретнее - перепишу)% TODO: хуйня какая то
\end{proof}
\begin{proposition}
	Если $f\in\mc R(E)$, то $f\in\mc L(E)$.
\end{proposition}
\begin{definition}
	Пусть $E$ - измеримое по Лебегу множество, а $f:E\ra\R$ ограничена на $E$. Положим:
	\[m=\inf_{x\in E}f(x),\ M=\sup_{x\in E}f(x),\ Q:m=q_0<\dots<y_N=M\]
	Назовём \textit{разбиением Лебега} следующие множества:
	\[E_i=\{x\in E:f(x)\in[y_{i-1},y_i),\ i<N\},\ E_N=\{x\in E:f(x)\in[y_{N-1},y_N]\}\]
	Назовём интегралом Лебега слелующий предел:
	\begin{multline*}
		\lim_{\tr(Q)\ra0}\sum_{i=1}^Nf(t_i)\mu(E_i)=\int_Ef(x)d\mu(x)\lra\\
		\lra(\fa\eps>0)(\ex\delta>0)(\fa Q,\ \tr(Q)<\delta)(\fa t_i,\ t_i\in E_i)\l|\sum_{i=1}^Nf(t_i)\cdot\mu(E_i)-\int_Ef(x)d\mu(x)\r|<\eps
	\end{multline*}
	Заметим, что, в отличие от интеграла Римана, в интеграле Лебега
\end{definition}
\begin{theorem}\nf{(Основная теорема об интегралах Лебега от огрниченных измеримых функций)}\label{11th1}
	Если $E$ - измеримое по Лебегу подмножество $\rn$ конечной меры, то любая ограниченная измеримая на $E$ функция $f$ интегрируема по Лебегу, причём её интеграл $\ds\int_Ef(X)d\mu(x)$ равен пределу интегральных сумм.
\end{theorem}
\begin{proof}
	Пусть $\mu(E)=0$. Доказывать нечего.\\
	Пусть $\mu(E)>0$. Тогда $\fa\eps>0$ положим $\ds\delta:=\frac\eps{\mu(E)}$, возьмём $\fa Q,\ \tr(Q)<\delta$.
	\[U(Q,f)-L(Q,f)=\sum_{i=1}^N(M_i-m_i)\cdot\mu(E_i)\le(y_i-y_{i-1})\cdot\mu(E_i)\le\tr(Q)\cdot\mu(E)<\eps\]
	Так как $\ds L(Q,f)\le\sum_{i=1}^Nf(t_i)\cdot\mu(E_i)\le U(Q,f)$, то требуемое получено.
\end{proof}
\begin{example}
	Функция Дирихле. Докажем, что $(\fa a,b,\ a\le b)\ f\in\mc L[a,b]$.
	\[
	\mc D(x)=
	\begin{cases}
		1,\ x\in\Q\\
		0,\ x\notin\Q
	\end{cases},
	\]
	Разобьём область значений $\mc D(x)$:
	\[[m,M]=[0,1],\ Q:0=y_0<\dots<y_q=1\]
	Пусть $\tr(Q)<1$, тогда:
	\[E_1=[a,b]\setminus\Q,\ E_q=[a,b]\cap\Q,\ \mu(E_1)=b-a,\ \mu(E_q)=0\]
	В итоге имеем:
	\[\lim_{\tr(Q)\ra0}\sum_{i=1}^qf(t_i)\mu(E_i)=f(t_1)\cdot(b-a)+f(t_q)\cdot0=0\]
	Получили требуемое.	
\end{example}
\begin{note}
	Интегрируемость по Лебегу была определена только для ограниченных измеримых функций. Теперь доопределим её у неотрицательных измеримых.
\end{note}
\begin{definition}
	Пусть $f$ - неотрицательная измеримая на $E$ функция, причём $\mu(E)<\i$, а у $E$ есть разбиение $P:E=\bsc_{i=0}^\i E_i$. Можно заметить, что появилось $E_0$, определим его:
	\[E_0=\{x\in E:f(x)=\p\}\]
	Если $\mu(E_0)>0$:
	\[U(P,f)=L(P,f)=\p\]
	Если $\mu(E_0)=0$, то:
	\[M_0\cdot\mu(E_0)=m_0\cdot\mu(E_0)=0\]
	Теперь снова определим суммы Дарбу и верхний и нижний интегралы:
	\[U(P,f)=\sum_{i=0}^\i M_i\cdot\mu(E_i),\ L(P,f)=\sum_{i=0}^\i m_i\cdot\mu(E_i),\ \v I=\inf U(P,f),\ \uv I=\sup L(P,f)\]
	Тогда разбиение Лебега:
	\[Q:0=y_0<y_1<\dots,\ \tr(Q)=\sup_{i\in\N}\tr y_i,\ E_i=\{x\in E:f(x)\in[y_{i-1},y_i]\}\]
\end{definition}
\begin{theorem}\nf{(Основная теорема об интеграле Лебега от неотрицательных измеримых функциях)}\label{11th2}\\
	Если $E$ - измеримое по Лебегу подмножество $\rn$ конечной меры, то любая неотрицательная измеримая на $E$ функция $f$ измерима по Лебегу, причём: \[\int_Ef(x)d\mu(x)=\lim_{\tr(Q)\ra0}\sum_{i=0}^\i f(t_i)\cdot\mu(E_i)\]
	где $Q$ - разбиение полупрямой $[0,\p]$.
\end{theorem}
\begin{proof}
	Аналогично \nr{11th1}. В суммах верхний индекс не $N$, а $\i$.\\
	Если $U(Q,f)=L(Q,f)=\p$, то:
	\[\int_Ef(x)d\mu(x)=\p\]
\end{proof}
\begin{definition}
	Если интеграл Лебега $\ds\int_Ef(x)d\mu(x)$ конечен, то $f$ называется \textit{суммируемой на $E$}.
\end{definition}
\begin{note}
	Теперь доопределим интеграл Лебега для любых измеримых функций.
\end{note}
\begin{definition}
	Пусть $f$ измерима на $E$, $\mu(E)<\p$. Построим 2 других функции:
	\[f^+(x):=\max(0,\ f(x)),\ f^-(x)=\max(0,\text{ -}f(x)\]
	Заметим, что и $f^+(x)$, и $f^-(x)$ являются измеримыми и неотрицательными функциями, следовательно, по \nr{1th2}, $f^+(x)\in\mc L(E),\ f^-(x)\in\mc L(E)$.\\
	Если хотя бы 1 из этих 2 интегралов Лебега конечен, то $f$ интегрируема по Лебегу, причём:
	\[\int_Ef(x)d\mu(x):=\int_Ef^+(x)d\mu(x)-\int_Ef^-(x)d\mu(x)\]
\end{definition}

\subsection{Основные свойства интеграла Лебега}
\begin{theorem}\nf{(Линейность и монотонность интеграла Лебега)}\label{12th1}
	\begin{enumerate}
		\item Если $f$ и $g$ суммируемы на $E$, $\mu(E)<\i$, то $(\fa \al,\beta\in\R)\ \al f+\beta g\text{ суммируема на }E$, причём:
		\[\int_E\al f+\beta g)(x)d\mu(x)=\al\int_Ef(x)d\mu(x)+\beta\int_Eg(x)d\mu(x)\]
		\item Если $\fa x\in E: f(x)\le g(x)$, $f$ и $g$ суммируемы на $E$, $\mu(E)<\i$, то:
		\[\int_Ef(x)d\mu(x)\le\int_Eg(x)d\mu(x)\]
	\end{enumerate}
\end{theorem}
\begin{proof}
	Пусть $\al=\beta=1$, а $f$ и $g$ ограничены. Заметим:
	\[L(P,f)+L(P,g)\le L(P,f+g)\le U(P,f+g)\le U(P,f)+U(P,g)\]
	Так как $f$ и $g$ суммируемы, то:
	\[(U(P,f)+U(P,g))-(L(P,f)+L(P,g))<\eps\Ra U(P,f+g)-L(P,f+g)<\eps\]
	Получили, что $f+g\in\mc L(E)$.\\
	Положим $c\ge0$, тогда:
	\[cU(P,f)=U(P,cf),\ cL(P,f)=L(P,cf)\]
	Если $c<0$, то:
	\[cU(p,f)=L(P,cf),\ cL(P,f)=U(P,cf)\]
	Пусть $f$ и $g$ неотрицательны и измеримы, следовательно:
	\[\mu(E_0(f))=0,\ \mu(E_0(g))=0\]
	Аналогично с ограниченным случаем.\\
	Пусть $f=f^+-f^-,\ g=g^+-g^-,\ f+g=(f+g)^+-(f+g)^-$, распишем:
	\[f+g=(f^++g^+)-(f^-+g^-)=(f^+-f^-)+(g^+-g^-)\]
	\[(f+g)^++f^-+g^-=(f+g)^-+f^++g^+\]
	\begin{multline*}
		\int_E(f+g)^+(x)d\mu(x)+\int_Ef^-(x)d\mu(x)+\int_Eg^-(x)d\mu(x)=\\
		=\int_E(f+g)^-(x)d\mu(x)+\int_Ef^+(x)d\mu(x)+\int_Ef^+(x)d\mu(x)
	\end{multline*}
	В итоге имеем:
	\begin{multline*}
		\int_E(f+g)^+(x)d\mu(x)-\int_E(f+g)^-(x)d\mu(x)=\\
		=\int_Ef^+(x)d\mu(x)+\int_Eg^+(x)d\mu(x)-\int_Ef^-(x)d\mu(x)-\int_Eg^-(x)d\mu(x)
	\end{multline*}
	Если внимательно посмотреть, то требуемое получено.
\end{proof}
\begin{definition}
	Определим \textit{срез неотрицательной функции}:
	\[
	f_{[N]}(x):=
	\begin{cases}
		f(x),\ 0\le f(x)\le N\\
		N,\ f(x)>N
	\end{cases}
	\]
\end{definition}
\begin{theorem}\nf{(Интеграл Лебега как предел интегралов от срезов)}\label{12th2}\\
	Если $f$ - неотрицательная измеримая на $E\subset\rn$ функция, $\mu(E)<\p$, то:
	\[\int_Ef(x)d\mu(x)=\lim_{N\ra\i}\int_Ef_{[N]}(x)d\mu(x)\]
\end{theorem}
\begin{proof}
	Потом допишу.
\end{proof}